\documentclass{situationist-preprint}

\title{The New Evil Demon Problem Answered}
\msid{EPI-06-b}
\msfamily{EPI-06}
\msdiscipline{Epistemology}
\mstype{Companion manuscript}
\mskeywords{new evil demon; internalism; externalism; reliabilism; excuse; brain in a vat}

\begin{document}
\maketitle

\begin{abstract}
The new evil demon problem asks how a reliabilist can count the demon's victim as justified. EPI-06-a answers that justification is internalist and knowledge externalist, so that the victim is justified and does not know. We test the answer against five variants: the envatted subject, the reliabilist's indexing of reliability to the actual world, the victim whose reasoning is careless, the victim who accepts testimony beyond her evidence, and the proposal that the victim is merely excused. In each case the verdicts follow from the supervenience of justification on internal states and the externalism of knowledge, and the rival treatments either deliver the same verdicts or concede internalism about the status at issue.
\end{abstract}

\section{The problem and the answer}

\citet{LehrerCohen1983} and \citet{Cohen1984} imagined subjects whose experiences and reasoning are exactly like ours but whose experiences are produced by a deceiving demon. Their beliefs are formed by processes that are unreliable in their world, yet they seem as justified as we are. If justification is reliable formation \citep{Goldman1979}, they are not.

EPI-06-a distinguishes two statuses. Justification answers the question whether a subject has responded as she should to what she has to go on (its Principle~P1), and is therefore internalist (Lemma~3.1); knowledge answers the question whether her belief can be relied on as a source of truth (Principle~P2), and is externalist (Lemma~3.2). By Corollary~3.4, the victim is justified to the same degree as her unvictimized duplicate and does not know. We examine the variants of the problem.

\section{Five variants}

\paragraph{The envatted subject.} A brain in a vat, envatted recently or for life, whose mental states and their history match those of an embodied perceiver, is her internal duplicate (Definition~2.1 of EPI-06-a) and so is justified in her perceptual beliefs to the same degree. That on semantic externalist grounds a lifelong envatted subject's beliefs have different contents (EPI-02-b) changes nothing: she responds as she should to what she has to go on, and her beliefs, whatever their contents, are justified. Nor is the verdict relative to the context of attribution (Corollary~3.5 of EPI-06-a): a sceptic who raises the possibility of envatment does not thereby make the envatted subject, or us, less justified, though by EPI-02-a and EPI-13 she can make it false to say that we know.

\paragraph{Indexed reliabilism.} The reliabilist may evaluate the victim's processes by their reliability in the actual world, or in normal worlds, rather than in hers \citep{Sosa1991,Comesana2002}. The victim is then justified, since her processes are of a type reliable here. But the proposal makes the status a function of the subject's process types alone, assessed at a fixed world, and internal duplicates have the same process types. The status so defined is shared by internal duplicates, and is internalist in the sense of Definition~2.2 of EPI-06-a. Indexed reliabilism thus answers the problem by conceding the thesis of EPI-06-a about justification, while locating the externalist component where EPI-06-a locates it, in knowledge.

\paragraph{The careless victim.} A victim who forms beliefs by hasty generalization, or who ignores evidence she has, is no more justified than her careless duplicate in our world. Internalism about justification does not make every victim justified; it makes victims exactly as justified as their duplicates, and the demon's interference neither lowers nor raises that standing.

\paragraph{The credulous hearer.} A victim who accepts the testimony of a demon-produced speaker beyond what her evidence about the speaker warrants is not justified, and neither is her duplicate who accepts the testimony of a real speaker on the same evidence. Since justification answers to what the hearer has to go on, a hearer whose acceptance outruns her evidence has not responded as she should, whatever the speaker's reliability. This is the internalist verdict that EPI-16-a relies on in its treatment of hearers who credit testimony beyond their evidence.

\paragraph{The excuse view.} Some externalists hold that the victim is not justified but blamelessly excused \citep{Littlejohn2012}. An excuse is appropriate when a subject fails a standard through no fault of her own. The victim fails no standard that answers to the concern of Principle~P1 of EPI-06-a: she has responded to her evidence exactly as we have. What she fails is the standard of Principle~P2, and failing that standard is failing to know, which EPI-06-a concedes. The excuse view relabels the internalist status as an excuse and reserves \enquote{justified} for a component of knowledge; the substantive verdicts are those of EPI-06-a.

\section{Objection}

\begin{objection}
If justification is internalist, there is no explanation of why justified beliefs tend to be true, and justification loses its connection to the aim of belief.
\end{objection}

\begin{reply}
The connection is secured by knowledge, of which justification is a component. A justified belief is one that the subject was right to form given what she had to go on; whether what she had to go on was a good guide to the truth is a further, external matter. In a world like ours it is, and so justified beliefs tend to be true here; in the demon's world it is not. That the tendency depends on the world is what the externalism of knowledge says.
\end{reply}

\section{Conclusion}

Each variant of the new evil demon problem is answered by the two theses of EPI-06-a. The victim is justified, because justification is shared by internal duplicates; she does not know, because knowledge is not. The rival treatments that succeed do so by adopting one or both of these theses.

\begin{thebibliography}{9}\small
\bibitem[Cohen(1984)]{Cohen1984} Cohen, S. (1984). Justification and truth. \emph{Philosophical Studies}, 46(3), 279--295.
\bibitem[Comesa{\~n}a(2002)]{Comesana2002} Comesa{\~n}a, J. (2002). The diagonal and the demon. \emph{Philosophical Studies}, 110(3), 249--266.
\bibitem[Goldman(1979)]{Goldman1979} Goldman, A.~I. (1979). What is justified belief? In G.~S. Pappas (Ed.), \emph{Justification and Knowledge} (pp.~1--23). Reidel.
\bibitem[Lehrer and Cohen(1983)]{LehrerCohen1983} Lehrer, K. and Cohen, S. (1983). Justification, truth, and coherence. \emph{Synthese}, 55(2), 191--207.
\bibitem[Littlejohn(2012)]{Littlejohn2012} Littlejohn, C. (2012). \emph{Justification and the Truth-Connection}. Cambridge University Press.
\bibitem[Sosa(1991)]{Sosa1991} Sosa, E. (1991). \emph{Knowledge in Perspective: Selected Essays in Epistemology}. Cambridge University Press.
\end{thebibliography}

\end{document}
